\documentclass[a4paper,10pt]{article}
\usepackage[utf8]{inputenc}
\usepackage[T1]{fontenc}
\usepackage[french]{babel}
\usepackage{amsmath,amssymb}
\usepackage{geometry}
\usepackage{array,tabularx}
\usepackage{tikz}
\geometry{margin=1.6cm}
\pagestyle{empty}
\newcommand{\ligne}{\par\vspace{0.55cm}\noindent\dotfill\par}
\newcommand{\carreaux}[1]{\begin{center}\begin{tikzpicture}\draw[step=0.5cm,gray!35,very thin] (0,0) grid (17,#1);\end{tikzpicture}\end{center}}
\newenvironment{exercice}[2]{\paragraph{Exercice #1 \hfill #2 points}\mbox{}\par}{\medskip}
\begin{document}
\begin{center}\Large\bfseries Devoir surveillé 16 - Les probabilités\end{center}
\vspace{2mm}
\begin{exercice}{1}{4}On lance un dé équilibré à six faces. Donner la probabilité d'obtenir $4$, puis un nombre pair.\end{exercice}\begin{exercice}{2}{5}Une urne contient $5$ boules rouges, $3$ bleues et $2$ vertes. Calculer la probabilité de tirer une boule rouge, puis une boule qui n'est pas bleue.\end{exercice}\begin{exercice}{3}{5}Classer les événements : impossible, certain, peu probable, équiprobable.\begin{enumerate}\item Obtenir pile en lançant une pièce équilibrée.\item Tirer une boule jaune dans une urne sans boule jaune.\item Obtenir un nombre inférieur à $7$ avec un dé à six faces.\end{enumerate}\end{exercice}\begin{exercice}{4}{6}Inventer une expérience aléatoire dans laquelle la probabilité d'un événement est $\dfrac14$, puis expliquer.\end{exercice}\end{document}
